%=============================================================================!
% 2.6  EQUATIONS OF MOTION AND THE STATE                                      !
%=============================================================================!
\section{Equations of motion and the state}
\label{sec:ne-equations}

For the falling ball the force was constant, and two integrations gave
the motion. Most forces are not constant: the drag of a liquid grows with
the speed of the body moving through it, and the pull of a spring grows
with how far it is stretched. This section shows what the second law
becomes in general, and what must be known at one instant to predict the
whole motion.

\subsection*{The equation of motion}

For motion along a line the force may depend on the position $x$, on the
velocity $v$ and on the time $t$, and we write it $F(x, v, t)$. Since the
acceleration is the second derivative of the position, the second law
reads
\begin{equation}
   m\,\frac{\dd^2 x}{\dd t^2} = F(x, v, t) .
   \label{eq:ne-eom}
\end{equation}
An ordinary equation such as $x^2 = 4$ has numbers as its unknowns;
\cref{eq:ne-eom} has a whole function, $x(t)$, as its unknown, and it
contains that function through its derivatives. An equation of this kind
is called a differential equation. Its order is that of the highest
derivative it contains, so \cref{eq:ne-eom} is of second order. It is the
equation of motion of the body, and solving it means finding a function
$x(t)$ for which it holds at every instant.

\subsection*{The state}

\Cref{eq:ne-eom} can be split into two first-order equations by treating
the velocity as a second unknown. The velocity is the rate of change of
the position, and the acceleration, which is the rate of change of the
velocity, equals $F/m$ by the second law, so
\begin{equation}
   \frac{\dd x}{\dd t} = v, \qquad
   \frac{\dd v}{\dd t} = \frac{F(x, v, t)}{m} .
   \label{eq:ne-first-order}
\end{equation}
These equations say that if the position and the velocity are known at
some instant, then the rates at which both are changing are known too.
Over a short interval $\tau$, therefore, the position changes by about
$v\tau$ and the velocity by about $(F/m)\tau$, which gives the position
and velocity at the time $t + \tau$, from which the next step follows in
the same way.

Take the thrown ball of Chapter~1, leaving the hand at $x = 0$ with
$v_0 = \SI{12}{\metre\per\second}$, where $F/m = -g$, and take one step of
$\tau = \SI{0.1}{\second}$:
\begin{align*}
   x(0.1) &\approx x(0) + v(0)\,\tau
      = 0 + 12\times0.1 = \SI{1.200}{\metre} ,\\
   v(0.1) &\approx v(0) - g\tau
      = 12 - 9.80665\times0.1 = \SI{11.019}{\metre\per\second} .
\end{align*}
The exact height from \cref{eq:mo-ball} is $12\times0.1 - \tfrac12
\times9.80665\times0.1^2 = \SI{1.151}{\metre}$. The step overshoots by
$\tfrac12 g\tau^2 = \SI{0.049}{\metre}$, because it used the velocity at
the start of the interval throughout, while the ball was in fact slowing
down. Halving $\tau$ divides this error per step by four but doubles the
number of steps needed to cover a given time, so the error at a given
time is halved. Shorter steps therefore approach the true motion.
Chapter~9 builds on this approximation to derive the integrators used in
molecular dynamics.

Once the force law $F(x, v, t)$ is given, the stepping argument
suggests, and the theorem below guarantees, that the pair $(x, v)$ at a
known instant is everything that needs to be known to predict the
future, and it is called the state of the body. The instant matters only
when the force depends on time explicitly, that is, through its third
argument $t$ and not only through the changing $x$ and $v$, as a push
switched on at a chosen moment does; the forces between atoms in this book do not, and
for them the state alone suffices.
Neither half is enough on its own. \Cref{fig:ne-state}(a) shows three
balls that start from the same height with different velocities and
follow three different motions; balls thrown at the same velocity from
different heights would differ as well. \Cref{fig:ne-state}(b) draws the
same three motions as paths in the plane whose two coordinates are $x$
and $v$, so that each point is one state. The arrows show, for every
state, the direction in which \cref{eq:ne-first-order} moves it, and each
path follows the arrows.

\begin{figure}[htb]
\centering
\includegraphics{ch02_newton/state}
\caption{The state of a falling ball. (a) Three balls released
\SI{3}{\metre} above the ground with velocities of $+2$, 0 and
\SI{-2}{\metre\per\second}: the same position with different velocities
gives different motions. (b) The same motions as paths through the
states $(x, v)$. Each grey arrow shows how the state at its tail changes
over \SI{0.03}{\second}, and dots mark the starting states.}
\label{fig:ne-state}
\end{figure}

That each starting state leads to exactly one motion is a theorem of
mathematics, the existence and uniqueness theorem for differential
equations, known as the Picard--Lindel\"of theorem. It holds provided
that $F$ is continuous in time and its slopes with respect to $x$ and
$v$ stay bounded in a neighbourhood of the starting state. It then
guarantees a single solution for at least some interval of time. Smooth
force laws satisfy this condition away from singularities such as two
atoms occupying the same position. A discontinuous cutoff requires
separate care, which is one reason to examine how forces are truncated
in Chapter~8. In three dimensions the same holds
with $\vb{r}$ and $\vb{v}$ in place of $x$ and $v$, and the state is the
pair of vectors $(\vb{r}, \vb{v})$. This is the state of a particle; a
body that can also turn needs, in addition, its orientation and how fast
it is turning (Chapter~5).

\begin{keyidea}
The equation of motion, $m\,\dd^2\vb{r}/\dd t^2 = \vb{F}$, is a
second-order differential equation. Given the force law and the starting
time, the position and velocity together, the state, determine the
motion that follows.
\end{keyidea}

\begin{selfcheck}
Two stones are dropped from rest, one from a bridge and one from a
window \SI{1}{\metre} lower. Do they have the same state at the moment of
release? Do they follow the same motion?
\end{selfcheck}
